\documentclass[12pt,a4paper]{article}
\usepackage[utf8]{inputenc}
\usepackage[vietnamese]{babel}
\usepackage{amsmath,amssymb,amsthm}
\usepackage{geometry}
\usepackage{enumitem}
\usepackage[most]{tcolorbox}
\usepackage{hyperref}
\usepackage{fancyhdr}
\usepackage{tikz}
\usepackage{eso-pic}
\usepackage{xcolor}
\geometry{a4paper, margin=2cm, bottom=2.5cm}
\hypersetup{colorlinks=true, linkcolor=blue!70!black}
\setlist[itemize]{topsep=4pt, itemsep=2pt}
\setlist[enumerate]{topsep=4pt, itemsep=2pt}
\AddToShipoutPictureFG{%
  \AtPageCenter{%
    \begin{tikzpicture}[remember picture, overlay]
      \node[rotate=45, scale=1, opacity=0.06]{\fontsize{60}{72}\selectfont\bfseries\sffamily Đặng Tấn Quí};
    \end{tikzpicture}%
  }%
}
\pagestyle{fancy}
\fancyhf{}
\fancyhead[L]{\small\textit{LÝ THUYẾT NHÓM}}
\fancyhead[R]{\small\textit{Đặng Tấn Quí}}
\fancyfoot[C]{\thepage}
\fancyfoot[L]{\footnotesize\textcopyright\ Đặng Tấn Quí}
\fancyfoot[R]{\footnotesize SĐT: 0377\,122\,210}
\renewcommand{\headrulewidth}{0.4pt}
\renewcommand{\footrulewidth}{0.4pt}
\fancypagestyle{plain}{\fancyhf{}\fancyfoot[C]{\thepage}\fancyfoot[L]{\footnotesize\textcopyright\ Đặng Tấn Quí}\fancyfoot[R]{\footnotesize SĐT: 0377\,122\,210}\renewcommand{\headrulewidth}{0pt}\renewcommand{\footrulewidth}{0.4pt}}
\newtcolorbox{dinhnghia}[1][]{enhanced, breakable, colback=blue!3!white, colframe=blue!60!black, fonttitle=\bfseries, title={Định nghĩa}, #1}
\newtcolorbox{dinhly}[1][]{enhanced, breakable, colback=green!3!white, colframe=green!50!black, fonttitle=\bfseries, title={Định lý}, #1}
\newtcolorbox{tinhchat}[1][]{enhanced, breakable, colback=yellow!5!white, colframe=orange!50!black, fonttitle=\bfseries, title={Tính chất}, #1}
\newtcolorbox{phuongphap}[1][]{enhanced, breakable, colback=gray!5!white, colframe=gray!50!black, fonttitle=\bfseries, title={Phương pháp}, #1}
\newtcolorbox{luuy}[1][]{enhanced, breakable, colback=red!3!white, colframe=red!50!black, fonttitle=\bfseries, title={Lưu ý}, #1}
\newtcolorbox{congthuc}[1][]{enhanced, breakable, colback=cyan!3!white, colframe=cyan!50!black, fonttitle=\bfseries, title={Công thức}, #1}
\newtcolorbox{vidu}[1][]{enhanced, breakable, colback=violet!3!white, colframe=violet!50!black, fonttitle=\bfseries, title={Ví dụ}, #1}

\begin{document}
\thispagestyle{empty}
\begin{tikzpicture}[remember picture, overlay]
  \fill[blue!80!black] (current page.south west) rectangle (current page.north east);
  \node[anchor=west, white, font=\fontsize{26}{32}\bfseries\selectfont]
    at ([xshift=3cm, yshift=-7.5cm]current page.north west) {LÝ THUYẾT NHÓM};
  \node[anchor=west, white, font=\fontsize{18}{24}\selectfont]
    at ([xshift=3cm, yshift=-9cm]current page.north west) {Nhóm hữu hạn, đồng cấu, nhóm đối xứng --- Năm 4};
  \node[anchor=west, white, font=\Large\bfseries] at ([xshift=3cm, yshift=-13cm]current page.north west) {Biên soạn:};
  \node[anchor=west, yellow!80!white, font=\fontsize{22}{28}\bfseries\selectfont] at ([xshift=3cm, yshift=-14.5cm]current page.north west) {ĐẶNG TẤN QUÍ};
  \node[anchor=south east, white, font=\Large\bfseries] at ([xshift=-2cm, yshift=3cm]current page.south east) {\the\year};
\end{tikzpicture}
\newpage
\tableofcontents
\newpage

%% ================================================================
%%  CHƯƠNG 1: NHÓM VÀ NHÓM CON
%% ================================================================
\section{Nhóm và nhóm con}

\subsection{Ôn tập nhóm}

\begin{dinhnghia}
  \textbf{Nhóm} $(G, \cdot)$: tập $G$ với phép toán hai ngôi thỏa kết hợp, có phần tử trung hòa $e$, mọi phần tử có nghịch đảo. \textbf{Nhóm Abel} (giao hoán): $ab = ba$ với mọi $a, b \in G$. \textbf{Cấp} của nhóm $|G|$: số phần tử (nếu hữu hạn).
\end{dinhnghia}

\begin{dinhnghia}
  \textbf{Cấp} của phần tử $a \in G$: số nguyên dương nhỏ nhất $n$ sao cho $a^n = e$ (nếu tồn tại); ký hiệu $o(a)$. Nếu không tồn tại thì $o(a) = \infty$.
\end{dinhnghia}

\subsection{Nhóm con}

\begin{dinhnghia}
  Tập con $H \subset G$ gọi là \textbf{nhóm con} của $G$ nếu $H$ đóng với phép toán và nghịch đảo: $e \in H$; $a, b \in H$ $\Rightarrow$ $ab \in H$; $a \in H$ $\Rightarrow$ $a^{-1} \in H$. Tương đương: $H \ne \varnothing$ và $a, b \in H$ $\Rightarrow$ $ab^{-1} \in H$.
\end{dinhnghia}

\begin{vidu}
  $\{e\}$ và $G$ là nhóm con tầm thường. $n\mathbb{Z} \subset \mathbb{Z}$. $SL_n(\mathbb{R}) = \{A \mid \det A = 1\}$ là nhóm con của $GL_n(\mathbb{R})$.
\end{vidu}

\begin{dinhnghia}
  \textbf{Nhóm con sinh bởi} $S \subset G$: $\langle S \rangle = \bigcap \{H \mid H \text{ nhóm con, } S \subset H\}$ = tập mọi tích hữu hạn các phần tử của $S$ và nghịch đảo của chúng.
\end{dinhnghia}

\subsection{Nhóm cyclic}

\begin{dinhnghia}
  Nhóm $G$ gọi là \textbf{cyclic} nếu tồn tại $a \in G$ sao cho $G = \langle a \rangle = \{a^n \mid n \in \mathbb{Z}\}$. Phần tử $a$ gọi là \textbf{phần tử sinh}.
\end{dinhnghia}

\begin{dinhly}
  Mọi nhóm cyclic vô hạn đẳng cấu với $\mathbb{Z}$. Mọi nhóm cyclic hữu hạn cấp $n$ đẳng cấu với $\mathbb{Z}/n\mathbb{Z}$.
\end{dinhly}

\begin{dinhly}
  Nhóm con của nhóm cyclic là cyclic. Nếu $G = \langle a \rangle$ cấp $n$ thì với mỗi ước $d$ của $n$ có đúng một nhóm con cấp $d$ là $\langle a^{n/d} \rangle$.
\end{dinhly}

\subsection{Định lý Lagrange}

\begin{dinhly}[title={Định lý Lagrange}]
  Cho $G$ nhóm hữu hạn, $H$ nhóm con. Khi đó $|H|$ chia hết $|G|$. Số $[G:H] = |G|/|H|$ gọi là \textbf{chỉ số} của $H$ trong $G$.
\end{dinhly}

\begin{dinhly}
  Cấp của mọi phần tử $a \in G$ chia hết $|G|$. Hệ quả: $a^{|G|} = e$ với mọi $a \in G$.
\end{dinhly}

\begin{luuy}
  Điều ngược lại của định lý Lagrange không đúng: nếu $d \mid |G|$ thì chưa chắc tồn tại nhóm con cấp $d$ (định lý Sylow cho trường hợp $d$ là lũy thừa số nguyên tố).
\end{luuy}

%% ================================================================
%%  CHƯƠNG 2: ĐỒNG CẤU NHÓM
%% ================================================================
\section{Đồng cấu nhóm}

\begin{dinhnghia}
  Ánh xạ $\varphi: G \to H$ giữa hai nhóm gọi là \textbf{đồng cấu} nếu $\varphi(ab) = \varphi(a)\varphi(b)$ với mọi $a, b \in G$.
\end{dinhnghia}

\begin{tinhchat}
  $\varphi(e_G) = e_H$, $\varphi(a^{-1}) = \varphi(a)^{-1}$. Ảnh nhóm con là nhóm con; nghịch ảnh nhóm con là nhóm con.
\end{tinhchat}

\begin{dinhnghia}
  \textbf{Hạt nhân}: $\ker \varphi = \{a \in G \mid \varphi(a) = e_H\} = \varphi^{-1}(\{e_H\}$. \textbf{Ảnh}: $\mathrm{Im}\,\varphi = \varphi(G)$. Cả hai đều là nhóm con; $\ker \varphi$ còn là nhóm con chuẩn tắc.
\end{dinhnghia}

\begin{dinhnghia}
  Nhóm con $N \trianglelefteq G$ (\textbf{chuẩn tắc}) nếu $gNg^{-1} = N$ với mọi $g \in G$, tương đương $gN = Ng$ (lớp trái = lớp phải).
\end{dinhnghia}

\begin{dinhly}[title={Định lý đồng cấu thứ nhất}]
  Cho đồng cấu $\varphi: G \to H$. Khi đó $\ker \varphi \trianglelefteq G$ và $G/\ker \varphi \cong \mathrm{Im}\,\varphi$.
\end{dinhly}

\begin{dinhly}[title={Định lý đồng cấu thứ hai}]
  Cho $N \trianglelefteq G$, $H$ nhóm con. Khi đó $HN = NH$ là nhóm con, $N \cap H \trianglelefteq H$, và $HN/N \cong H/(N \cap H)$.
\end{dinhly}

\begin{dinhly}[title={Định lý đồng cấu thứ ba}]
  Cho $N \trianglelefteq G$, $K \trianglelefteq G$, $K \subset N$. Khi đó $N/K \trianglelefteq G/K$ và $(G/K)/(N/K) \cong G/N$.
\end{dinhly}

%% ================================================================
%%  CHƯƠNG 3: NHÓM ĐỐI XỨNG VÀ NHÓM THAY THẾ
%% ================================================================
\section{Nhóm đối xứng và nhóm thay thế}

\begin{dinhnghia}
  \textbf{Nhóm đối xứng} $S_n$: nhóm các song ánh $\sigma: \{1,2,\ldots,n\} \to \{1,2,\ldots,n\}$ với phép hợp. Phần tử $\sigma$ gọi là \textbf{phép thế} (hoán vị) bậc $n$. $|S_n| = n!$.
\end{dinhnghia}

\begin{dinhnghia}
  \textbf{Ký hiệu chu trình}: $(i_1 \, i_2 \, \cdots \, i_k)$ nghĩa là $i_1 \mapsto i_2 \mapsto \cdots \mapsto i_k \mapsto i_1$. Mọi phép thế phân tích được thành tích các chu trình rời nhau.
\end{dinhnghia}

\begin{dinhnghia}
  \textbf{Chu trình độ dài 2} gọi là \textbf{chuyển vị}. Mọi phép thế viết được thành tích các chuyển vị. \textbf{Dấu} của phép thế $\varepsilon(\sigma) = (-1)^{\text{số chuyển vị}}$ (chẵn/lẻ không phụ thuộc cách viết).
\end{dinhnghia}

\begin{dinhnghia}
  \textbf{Nhóm thay thế} (nhóm xen kẽ) $A_n = \{\sigma \in S_n \mid \varepsilon(\sigma) = 1\}$: nhóm các phép thế chẵn. $|A_n| = n!/2$. $A_n \trianglelefteq S_n$.
\end{dinhnghia}

\begin{dinhly}
  $A_n$ là nhóm đơn (không có nhóm con chuẩn tắc thực sự) khi $n \ge 5$. Đây là kết quả quan trọng trong chứng minh định lý Abel về phương trình bậc 5 không giải được bằng căn thức.
\end{dinhly}

\begin{vidu}
  $S_3$ có 6 phần tử, đẳng cấu với nhóm các phép đối xứng của tam giác đều. $S_3 = \langle (12), (123) \rangle$. $A_3 = \langle (123) \rangle \cong \mathbb{Z}/3\mathbb{Z}$.
\end{vidu}

%% ================================================================
%%  CHƯƠNG 4: TÍCH TRỰC TIẾP VÀ MỘT SỐ ĐỊNH LÝ
%% ================================================================
\section{Tích trực tiếp và định lý Cauchy}

\begin{dinhnghia}
  \textbf{Tích trực tiếp} (tích ngoài) $G \times H$: nhóm với phép toán $(g_1, h_1)(g_2, h_2) = (g_1 g_2, h_1 h_2)$. $|G \times H| = |G| \cdot |H|$.
\end{dinhnghia}

\begin{dinhly}[title={Định lý Cauchy}]
  Cho $G$ nhóm hữu hạn, $p$ là số nguyên tố chia hết $|G|$. Khi đó $G$ có phần tử cấp $p$ (và do đó có nhóm con cấp $p$).
\end{dinhly}

\begin{dinhly}[title={Định lý cấp nhóm cyclic}]
  Nhóm cyclic cấp $n$ có đúng $\varphi(n)$ phần tử sinh, với $\varphi$ là hàm Euler.
\end{dinhly}

\begin{dinhnghia}
  \textbf{Nhóm con chuẩn tắc} $N \trianglelefteq G$: $gN = Ng$ với mọi $g$. \textbf{Nhóm thương} $G/N$: tập các lớp $gN$ với phép nhân $(gN)(hN) = (gh)N$.
\end{dinhnghia}

\begin{dinhly}
  $N \trianglelefteq G$ khi và chỉ khi $N$ là hạt nhân của một đồng cấu từ $G$.
\end{dinhly}

\end{document}
